\chapter{Лазерное охлаждение в МОЛ}\label{ch:ch1}

\section{Введение}

Магнито-оптическая ловушка (МОЛ) работает засчёт комбинации лазерного охлаждения и неоднородного магнитного поля. 
В ловушке используется три пары противоположных лазерных пучков, которые создают силу трения, замедляющую и охлаждающую атомы, а магнитное поле с линейным градиентом — возвращающую силу, удерживающую атомы в центре ловушки. 
Лазерные лучи настраиваются с небольшой частотной отстройкой от резонанса, чтобы благодаря эффекту Доплера атомы поглощали фотоны, теряя кинетическую энергию. 
Магнитное поле индуцирует расщепление уровней атомов (эффект Зеемана), что помогает направлять атомы к центру ловушки из-за изменения частотно-зависимого взаимодействия с лазерным светом. 
Таким образом, МОЛ позволяет получить холодные, практически неподвижные нейтральные атомы при температурах порядка сотен или десятков микрокельвинов.

Доплеровская температура — это физический предел минимальной температуры, которой можно достичь при лазерном охлаждении атомов с использованием эффекта Доплера. 
Она связана с шириной энергетического перехода атома (шириной резонансной линии) и определяется такими параметрами, как частота и массовый состав атомов. 
По сути, это температура, при которой эффект лазерного охлаждения (основанный на разнице поглощения фотонов атомами, движущимися навстречу или по направлению лазера) перестает эффективно снижать кинетическую энергию атомов, потому что тепловое движение становится настолько малым, что атомы не различают лазерные потоки ни движения навстречу, ни в направлении движения атомов.

\begin{equation}
	\label{eq:DopplerLimit}
	T_D = \frac{h \Gamma}{2 k_B}
\end{equation}

Доплеровский предел задаётся формулой \ref{eq:DopplerLimit}
В числовом выражении доплеровская температура обычно находится на уровне десятков или сотен микрокельвин, что выше фундаментального предела отдачи (минимальной температуры, обусловленной отдачей импульса от излучаемых фотонов).
Причём чем уже охлаждающий переход, тем ниже итоговая температура, но тем меньше атомов удастся захватить. 
Поэтому, для получения низких температур в МОЛ используют комбинацию двух охлаждающих переходов.


В тулии есть много доступных для охлаждения переходов. 
Они приведены в таблице в прошлом пункте .
Был выбран переход на длине волны $\lambda = 506.2\unit{нм}$ со спектральной шириной $\gamma =  7.8\unit{кГц}$.
Это сравнимо с используемыми в текущий момент переходами для охлаждения других атомов в МОЛ.  
Для того, чтобы реализовать это охлаждение нужно сделать разное: узкополосный лазерный источник, знание частот переходов между различными сверхтонкими. 

\section{Источник излучение}

В качестве источника излучения использовался лазер Toptica DLpro. 
По спецификации, спектральная ширина данного лазера в режиме свободной генерации составляет около 300 кГц, при долговременной стабильности порядка 10 МГц/ч.  
Для работы с исследуемым переходом требуется стабилизировать частоту данного лазера и сузить его спектр. 

Стандартным методом сделать это является стабилизация частоты лазера по моде внешнего высокодобротного резонатора методом Паунда-Древера-Холла.
Основная идея метода заключается в том, что лазерный луч модулируется по фазе при помощи модулятора (обычно электрооптического), а затем направляется на высокодобротный оптический резонатор Фабри-Перо. 
Излучение, отраженное от резонатора, несет информацию об отстройке частоты лазера от частоты собственной моды резонатора.

Дальнейшая гетеродинная обработка отраженного сигнала позволяет получить так называемый сигнал ошибки — величину, пропорциональную разности частоты лазера и резонансной частоты резонатора. 
Этот сигнал ошибки используется для обратной связи и точной подстройки частоты лазера, что обеспечивает его высокую стабильность и узкую спектральную линию.

В нашем случае ток лазерного диода модулировался с помощью внутренних средств блока управления лазером. 
В качестве опорного использовался резонатор из УЛЕ стекла, зеркала устанавливались на него с помощью оптического контакта.
Конструкция вакуума и температурных экранов подробно описана в других работах.  
Итоговая резкость собранного резонатора составляла $\mathcal{F} = 47.5\times10^3$, при пропускании $T = 30\%$.
Отражённое от резонатора излучение подаётся на фотодиод \textbf{ThorlabsPhDModel}, полученный сигнал подавался на модуль PDH блока управления лазером, который генерировал сигнал ошибки.
Сервосигнал получался из него с помощью модуля Toptica FALC, и подавался на медленный токовый вход лазера. 
Медленная петля обратной связи формировалась цифровым образом и подмешивала напряжение к напряжению на пьезе лазера. 
Оптимизация FALC проводилась по Фурье-образу сигнала ошибки, а также по спектру лазерного излучения, измеряемого в эксперименте. 

Такая стабилизация позволяла сузить линию лазера до величины лучше 1 кГц (особенности спектра такой ширины детектировались в эксперименте).
Измерение ширины линии излучения методом автогоммодинирования с помощью волоконной линии задержки не проводилось: для этого требовалась волоконная линия длинной более 300 км, что для волокна на требуемую длину волны невозможно много.

\section{Спектроскопия перехода}

Для получения магнито-оптической ловушки требовалось определить частоту нужного сверхтонкого, а значит нужно провести спектроскопию данного перехода.
При наличии магнито-оптической ловушки можно осуществить спектроскопию потерь атомов из МОЛ. 
При работе первой стадии охлаждения возбуждается переход $\ket{g, F=4}\xrightarrow{} \ket{b, F = 5}$, при этом происходит нерезонансное возбуждение перехода $\ket{g, F=4}\xrightarrow{} \ket{b, F=5}$, и часть атомов выбывают из цикла охлаждения. 
Излучение на исследуемой длине волны присутствует в области захвата атомов постоянно, его частота сканируется. 
При попадании частоты излучения в резонанс с переходом из основного состояния часть атомов выходит из цикла охлаждения, что приводит к потерям из МОЛ.
Если же излучение настроенно в резонанс с верхнего сверхтонкого, то атомы, нерезонансно возбуждённые ранее, возвращались в цикл охлаждения, что приводило к увеличению числа атомов в итоге. 

В данном эксперименте частота излучения на длине волны \thd не стабилизировалась.
Перестройка частоты излучения осуществлялась изменением напряжения на пьезокерамике, а длина волны измерялась по измерителю длин волн \textbf{WLMModelName}. 

%Полученный спектр представлен на рисунке \ref{fig:HFS_spectrum}, частоты найденных переходов приведены в таблице \ref{tab:Hyperfine506}.
Измеренные частоты переходов между различными сверхтонкими компонентами основного и возбуждённого состояний исследуемого перехода приведены в таблице \ref{tab:Hyperfine506}.
Частота перехода $F = 3 \xrightarrow{} F = 4$ определяется однозначно, т.к. это единственный разрешённый переход из состояния \gF{3}. 
Определение того, какой из двух переходов c состояния $F = 4$ какой, происходило на основании того, что характерные для атомов тулия сверхтонкие расщепления отрицательны, поэтому мы решили, что $F = 4 \xrightarrow{} F = 4$ переход должен обладать большей частотой.
Это подтверждалось также тем, что в последствии на $F = 4 \xrightarrow{} F = 5$ переходе удалось получить МОЛ.
Таким образом, измерено также сверхтонкое расщепление основного и возбуждённого состояний охлаждающего перехода. 
Для основного состояния это расщепление составило $^{HFS}f = -1497(7)\unit{МГц}$, что согласуется с измеренным в работе \cite{mishin2024combined} значением $^{HFS}f = - 1 496 550 658.23(3)\unit{Гц}$. 
Для возбуждённого же состояния расщепление составило $^{HFS}f = -3474(4)\unit{МГц}$, что согласуется со значением $^{HFS}f = -3474(2)\unit{МГц}$ из работы \cite{kroger1997high}. 
%\TODO{Проверить статью, там точно была какая-то статья Сорокина}

\begin{table}[]
	\centering
	\begin{tabular}{|cc|}
		\hline
		Переход между сверхтонкими компонентами & Частота (ГГц)\\
		\hline
		$F = 4 \xrightarrow{} F = 5$		&  592045.689\\
		$F = 4 \xrightarrow{} F = 4$		&  592049.163\\
		$F = 3 \xrightarrow{} F = 4$ 		&  592047.666\\
		$F = 3 \xrightarrow{} F = 5$ 		&  Запрещён\\
		\hline
	\end{tabular}
	\caption{Частоты переходов между различными компонентами сверхтонкого расщепления уровней для третьей стадии лазерного охлаждения}
	\label{tab:Hyperfine506}
\end{table}


\section{МОЛ на длине волны \thd}
После того, как была определена частота охлаждающего перехода, излучение на длине волны \thd было совмещено с двумя другими охлаждающими излучением (\fst и \snd) перед волокном, доставляющим охлаждающее излучение к вакуумной камере.
Таким образом создавалась классическая 6-лучевая конфигурация МОЛ.
Характерный размер пучков охлаждающего излучения составлял $w_{1/e^2} = 1\unit{см}$.

После охлаждения во вторичной МОЛ на длине волны \snd температура атомов в МОЛ составляет $T\approx20\unit{мкК}$, что соответствует средней скорости атомов $v \approx 10\unit{см/с}$, которая приводит к Доплеровскому сдвигу излучения на длине волны \thd $\Delta\nu_D \approx 100\unit{кГц} \approx 12 \gamma$. 
Резонансная частота перехода определялась по спектру поглощения на длине волны \thd облака атомов, удерживающихся в оптической патоке на длине волны \snd. 
После охлаждения атомов в магнито-оптической ловушке, магнитные поля МОЛ выключались, и через $30\unit{мс}$ -- время достаточное для выключения магнитных полей, возбуждался переход, и измерялось количество оставшихся в МОЛ атомов.  
Пример полученного спектра приведён на рисунке \ref{fig:Line_center_MOT}
Структура с двумя пиками объясняется Зеемановскими сдвигами для различных магнитных подуровней основного состояния. 
Магнитное поле внутри вакуумной камеры может быть неоднородным, и, вследствие большого атомного облака в данном эксперименте, разные атомы могут находиться в разных магнитных полях и, возможно, на различных магнитных подуровнях.
Если принять во внимание $g$-факторы основного и возбужденного состояний для перехода $506.2 \unit{нм}$ \cite{martin1978atomic}, то данную структуру можно объяснить различием остаточного магнитного поля величиной приблизительно в $10 \unit{мГс}$ на масштабе диаметра атомного облака $\sim0.5\unit{мм}$.

\begin{figure}[ht]
	\centerfloat{
		\includegraphics[scale=0.75]{Chapter1/Line_center_MOT.pdf}
	}
	\caption{
		Типичный спектр поглощения на охлаждающем переходе $506.2 \unit{нм}$ в свободном пространстве в зависимости от частоты АОМа $\nu$, отстроенной от $\nu_0 = 200.985\unit{МГц}$.
		Пунктирной линией показаны результаты аппроксимации двумя гаусссианами.
		Полусумма частот их центров была выбрана в качестве начала отсчета шкалы частот в последующих экспериментах.
	}\label{fig:Line_center_MOT}
\end{figure} 

После отстройки охлаждающего излучения на $\approx12\gamma$ мы увидели слабый перезахват в МОЛ на длине волны \thd.
Следующая секция посвящена оптимизации перезахвата.


\subsection{Перезахват атомов в спектрально-уширенную МОЛ}

Особенность захвата атомов в спектрально-узкую МОЛ состоит в том, что спектральная ширина линии излучения в начале перезахвата оказывается гораздо меньше доплеровского уширения, из-за чего эффективно происходит взаимодействие лишь с малой скоростной группой атомов. 
Стандартным подходом, позволяющим преодолеть эту сложность является уширение линии излучения с помощью частотной модуляции. 
В нашем случае, для формирования частоты излучения использовался акусто-оптический модулятор, радиочастотный сигнал для которого генерировался платой DDS на основе чипа AD9910.
Данный чип позволяет осуществлять модуляцию частоты выходного сигнала треугольной формы.
Спектральные плотности одночастотного и модулированного сигнала представленны на рисунке \ref{fig:Modulation_Spectra}. 


\begin{figure}[ht]
	\centerfloat{
		\includegraphics[scale=0.6]{Chapter1/Modulation_Spectra.pdf}
	}
	\caption{
		Спектры сигналов, подаваемых на управляющий АОМ в одночастотном и модулированном режиме.
		Амплитуда модуляции частоты составляла $\Delta\nu = 100 \unit{кГц}$, частота составляла $\Omega=10\unit{кГц}$.
		Спектра получены с помощью спектроанализатора SpecrtoAnalyserModelName.   		
	}\label{fig:Modulation_Spectra}
\end{figure} 

Выбор амплитуды модуляции делался из следующих соображений.
При характерной температуре атомного облака $T = 20\unit{~мкК}$, средняя скорость атомов составляет $\sim10\unit{м/с}$, что соответствует Доплеровскому сдвигу излучения на длине волны \thd $\Delta\nu_D \approx 100\unit{кГц} \approx 12 \gamma$. 
Амплитуда модуляции $\Delta\nu = 100\unit{кГц}$ позволяла эффективно взаимодействовать со всеми скоростными группами атомов.
Частота модуляции выбралась около $10\text{~кГц}$, чтобы перекрывать весь диапазон частот за время сравнимое с временем жизни верхнего уровня.
Рабочая отстройка от резонанса (верхняя частота) выбиралась порядка $-100\text{~кГц} \approx -12 \gamma$, чтобы оставлять модуляционные сайдбэнды вдали от резонанса. 

Далее изучался перезахват в магнито-оптическую ловушку на длине волны $506.2\text{~нм}$ в зависимости от длительности охлаждения на длине волны $530.7\text{~нм}$.

При отсутствии второй стадии охлаждения, удавалось осуществить перезахват около 15\% атомов из магнито-оптической ловушки на длине волны \thd. 
Характерное время, за которое вторичное охлаждение уменьшает температуру и радиус облака составляет $\tau = 5 \unit{мс}$, причем на константу оно выходит за время порядка $\tau \sim 20\unit{мс}$.
При идеальном выставлении магнитного поля удаётся достичь полного перезахвата из вторичной МОЛ в третичную, при этом характерное значение составляет около 65\%.  

\begin{figure}[ht]
	\centerfloat{
		\includegraphics[scale=0.5]{Chapter1/Recapturing.pdf}
	}
	\caption{
		Зависимость доли перезахваченных в МОЛ атомов от длительности вторичного охлаждения.
		Видно, что за характерное время $\tau\approx 5\text{~ms}$ и температура, и радиус облака уменьшаются до некоторой постоянной величин. 
		С такой же константой времени нарастает и доля перезахваченных атомов вплоть до 100\%.
	}\label{fig:recapturing}
\end{figure}

\subsection{Охлаждение в спектрально-узкой МОЛ}

После стадии охлаждения в спектрально-уширенной МОЛ, модуляция отключалась, и выставлялось ещё меньшее значение градиента магнитного поля.
Для оптимизации охлаждения мы измеряли следующие параметры ансамбля: число захваченных атомов, размер облака и его фазовую плотность, а также температуру; в зависимости от: градиента магнитного поля, мощности охлаждающего излучения, а также его отстройки от резонанса. 

\subsection{Измерение параметров облака}

Число атомов в облаке определялась по сигналу флуоресценции. 
В этом методе, фотография облака атомов на калиброванной камере аппроксимировалась Гауссовой функцией, амплитуда которой соответствовала числу атомов. 
Зная чувствительность камера на нужной длине волны, телесный угол, из которого собиралось излучение и скорость рассеяния фотонов, сигнал с камеры возможно перевести в число атомов.

Температура измерялась методом баллистического разлёта.
В начальный момент времени все удерживающие поля (квадрупольное магнитное поле и охлаждающие пучки МОЛ) и облако свободно падает в гравитационном поле. 
В предположении Максвелловского распределения атомов по скоростям и Гауссовой формы облака зависимость радиуса облака от времени задаётся формулой \ref{eq:BallisticExcpansion}, где $w_0$ - начальный размер облака, $k_B$ - постоянная Больцмана, $T$ - температура облака атомов и $m$ - масса атомов.

\begin{equation}
	\label{eq:BallisticExcpansion}
	w_{1/e}(t) = \sqrt{w_0^2 + \frac{2k_B T}{m}t^2}
\end{equation}


Ожидаемая температура атомов (порядка пределу отдачи $T_{rec} \sim 200\unit{мкК}$) соответствует очень медленному разлёту ансамбля атомов.
Так, расширение облака вдвое, требовалось время разлёта порядка $tau \sim 30\unit{мс}$, что соответствовало падению на $4.5\unit{мм}$, т.е. ограниченно уже полем зрения камеры.
Поэтому, для того, чтобы пользоваться таким методом, излучение, используемое для регистрации атомов, заводилось в систему сверху, под небольшим углом к вертикали.
Из-за медленного разлёта требовалось большое время ожидания ($\tau \sim 30\unit{мкК}$) и нужно было использовать вс поле зрения камеры. 

Для верификации полученных результатов проводилась спектроскопия часового перехода в свободном пространстве. 
После охлаждения атомов в МОЛ, квадрупольное магнитное поле выключалось, и за время, необходимое для спада магнитного поля (порядка $\tau\approx10\unit{мс}$) атомы удерживались в оптической патоке на длине волны \thd. 
Затем, удерживающие оптические поля выключались и проводился опрос атомов вертикальным пучком излучения на длине волны \clck. 
Вертикальное направление выбиралось из-за того, что это наиболее интересное для часов направление, и у нас не было пучков в других направлениях. 
Ширина полученного контура описывалась в приближении только Доплеровского уширения, и полученная спектральная ширина использовалась как оценка сверху на температуру облака атомов.

Фазовая плотность ($PSD = n\Lambda_{dB}^3$, где $PSD$ - фазовая плотность, $n$ - концентрация атомов, а $\lambda_{dB}$ - длина волны де-Бройля) определялась из комбинации двух предыдущих измерений. 
Из фотографии облака в момент выключения удерживающих полей определялось число атомов в ансамбле и геометрический размер облака. 
Объём, занимаемый атомами в импульсном пространстве, определялся из измерения температуры ($\Lambda_{dB} = \frac{h}{\sqrt{2M k_B T}}$). 
Здесь стоит сказать, что температура атомов может быть различна в вертикальном и горизонтальных направлениях. 
Итоговая формула: $PSD = \frac{N}{\pi^{3/2} w_z w_r^2} \times \Lambda_z \Lambda_r^2$, где N - измеренное число атомов, $w_z$ и $w_r$ - радиусы облаков в вертикальном и радиальном направлениях, а $\Lambda_z$ и $\Lambda_r$ фазовые радиусы, определяемые формулой . 

\begin{figure}[ht]
	\centerfloat{
		\includegraphics[scale=0.35]{Chapter1/TemperatureMeasurement}
	}
	\caption{
		Два метода температуры облака атомов. 
		На рисунке б) показан стандартный метод измерения температуры методом баллистического разлёта. 
		Температура, полученная в результате аппроксимации полученных данных формулой \ref{eq:BallisticExcpansion} составила $T = 690(60)\unit{нК}$.
		Доплеровская ширина контура часового перехода соответствующая такой температуре составляет $\Delta \nu_D = 12.1\unit{кГц}$, что находится в согласии с измеренной шириной $\Delta\nu = 13.7(1.2)\unit{кГц}$. 
	}\label{fig:TemperatureMeasurement}
\end{figure}

\subsection{Оптимизация процесса охлаждения в МОЛ}

После отключения спектрального уширения охлаждающих пучков на длине волны \thd атомы охлаждались в спектрально-узкой МОЛ за время $\tau = 100\unit{мс}$. 
После охлаждения атомов проводились измерения параметров МОЛ методами описанными выше.
Для оптимизации охлаждения мы измеряли число захваченных атомов, температуру облака и его фазовую плотность в зависимости от: градиента магнитного поля, мощности охлаждающего излучения, и также его отстройки от резонанса.
Градиент магнитного поля изменялся в диапазоне $200-600\unit{мГс/см}$: уменьшение градиента приводило к слишком большому геометрическому размеру облака, и, как следствие, большим потерям атомов, а дальнейшее увеличение градиента не приводило к каким-либо изменениям.  
Отстройка охлаждающего излучения от резонанса менялась в диапазоне $\Delta\nu = -14 \div -8 \gamma$.
При приближении частоты излучения к резонансу резко уменьшалось число захваченных атомов, а увеличение отстройки приводили к слишком большому геометрическому размеру. 
Мощность излучения менялась в диапазоне $100 \div 500 \unit{мкВт}$, что соответствовало диапазону  $12 \div 90$ единиц интенсивности насыщения $s$.
Меньшая мощность приводила к уменьшению перезахвата, увеличение мощности не приводило к изменениям.

Результаты оптимизации представлены на рисунке \ref{fig:CoolingOptimization}. 
Они оказались ожидаемы: низкие температуры достигались при малом значении мощности, большие градиенты магнитного поля и малые отстройки приводили к получению компактной МОЛ с большим количеством атомов.
Минимальная достигнутая температура составляла $T=400\unit{нК}$, максимальная доля доживших до конца эксперимента атомов составляло 100\% ($N = 2\times10^6$).
Оптимальные значения фазовой плотности достигались компромиссом между температурой и размером облака.
Лучшая полученная фазовая плотность составляла $PSD = 3.5\times10^{-4}$.
С одной стороны, это гораздо меньше, чем желаемые 2.612, но уже на порядок выше, чем то, что достигали наши коллеги из РКЦ.
Если начинать с их числа атомов ($N_{RQC} = 30\times10^6$), можно достичь и двух порядков улучшения.
Некоторая часть результатов была получена дли МОЛ чечевичной формы, что приводило к большим погрешностям в измерении температуры. 
Такие области параметров обозначались белой штриховкой на графике.

\begin{figure}[ht]
	\centerfloat{
		\includegraphics[scale=0.15]{Chapter1/CoolingOptimization.pdf}
	}
	\caption{
		Результат оптимизации процесса охлаждения атомов в магнито-оптической ловушки для следующих параметров: градиент магнитного поля ($b$, ось $y$), отстройка охлаждающего излучения от резонанса ($\delta \nu$, ось $x$), и полная мощность  охлаждающего излучения на длине волны \thd в пучках МОЛ, в единицах интенсивности насыщения ($s$, строки).
		Четыре столбца показывают результат процесса оптимизации для четырёх исследуемых параметров: температуры атомного облака в вертикальном направлении  $T_\textrm{ver}$ и в горизонтальной плоскости $T_\textrm{hor}$, фазовая плотность ансамбля атомов ($PSD$) и относительное число атомов на момент окончания охлаждения ($N / N_0$, где $N_0 = 2\times 10^6$ --- число атомов в МОЛ после первой стадии охлаждения).
		Чечевичная форма облака атомов в области больших отстроек и малых градиентов магнитного поля приводит к большим ошибкам измерения температуры облака в горизонтальной плоскости (а значит и фазовой плотности), поэтому данные области параметров заштрихованы белыми линиями.  
	}\label{fig:CoolingOptimization}
\end{figure}

\section{Перезахват в несколько непересекающихся облаков либо в длинное облако}

В некоторых работах \cite{bothwell2022resolving,zheng2022differential} предлагается использовать несколько разных ансамблей атомов, захваченных в одну оптическую решётку, но разнесённых в пространстве на расстояние от несколько миллиметров до нескольких сантиметров. %обновить ссылки
Применяемый нами метод охлаждения позволяет получить несколько непересекающихся ансамблей атомов в оптической решётке засчёт большой чувствительности геометрического положения МОЛ к компенсационным магнитным полям.
Для этого ансамбль атомов тулия охлаждался в магнито-оптической ловушке на длине волны \snd в присутствии излучения оптической решётки, после чего излучение вторичной МОЛ отключается а градиентное магнитное поле переключается на режим третичного охлаждения.
Из-за того, что атомы в облаке уже достаточно холодные ($T\sim20\unit{мкК}$), часть атомов оказываются захвачены в оптическую решётку на месте вторичной МОЛ, при этом не захваченные в решётку атомы перезахватываются в МОЛ в координатах, задаваемых значением компенсационного магнитного поля.
После охлаждения атомов в спектрально-уширенной МОЛ на длине волны \thd отключались удерживающие поля, и часть атомов перезахватывалась в оптическую решётку на новом месте. 
Результаты применения этого протокола показаны на рисунке \ref{fig:DoubleCloud}.

\begin{figure}[ht]
	\centerfloat{
		\includegraphics[scale=0.75]{Chapter1/DoubleCloud.pdf}
	}
	\caption{
		Результат эксперимента по перезахвату двух независимых ансамблей МОЛ в различных координатах. 
		Число атомов в каждом из облаков получается около 7\% от начального числа атомов.  
	}\label{fig:DoubleCloud}
\end{figure}

%
%\section{Основные результаты главы 1}

Полученные в данной главе ансамбли атомов оказываются охлажденными до температур лишь вдвое больше предела отдачи исследуемого перехода. 
Такие ансамбли были получены ранее для различных атомов \cite{}, в этом смысле мы находимся на мировом уровне. 
Число перезахваченных атомов оказалось ниже чем для легких изотопов, и сравнимо с тяжёлыми атомами, что логично. 
Понимания того, какие достижимы фазовые плотности у меня пока нет, но надо будет это сюда заботать.

В полученных результатах есть несколько особенностей, которые нужно обсудить. 
Во-первых, после охлаждения в узкой и отстроенной МОЛ облако атомов оказывается большого размера, что резко ухудшает перезахват в оптическую решётку.
Коллеги из квантовых симуляторов умеют с этим справляться, собирая большую МОЛ в дипольные ловушки, в случае же часовых экспериментов с оптическими решётками это большой минус. 
При этом, использовать спектрально-уширенную МОЛ компактной формы для увеличения перезахвата атомов в решётку можно и нужно, удаётся добиться увеличения перезахвата больше, чем вдвое, но температура получается порядка $T\sim2\unit{мкК}$ что больше того, что нам необходимо. 
Во-вторых, при охлаждении в магнито-оптических ловушках с большой отстройкой от резонанса, распределение атомов по магнитным подуровням неравномерно, причём большая часть оказывается на крайних магнитных подуровнях. 
В том протоколе, который мы использовали ранее, большая часть потерь приходилась на атомы с крайних магнитных подуровнях. 
Всё это приводит к тому, что охлаждение в МОЛ на длине волны \thd хотя и улучшает протокол, но не решает наших проблем целиком, а значит необходимо разрабатывать также протокол подготовки начальных состояний часового  перехода в оптической решётке. 

\section{Основные результаты главы \ref{ch:ch1}}

\begin{itemize}
		\item При охлаждении нейтральных атомов тулия в магнито-оптической ловушке на переходе $\ket{4f^{13}(^2F^\text{o})6s^2, F = 4} \xrightarrow{} \ket{4f^{13}(^2F_{7/2}^\text{o})6s6p(^3P_2^\text{o}, F=5)}$ с длиной волны $\lambda = 506.2~\unit{нм}$ достижимы температуры облака атомов $T = 400\unit{нК}$, близкие к пределу отдачи используемого перехода ($T_r = 220\unit{нК}$), а также фазовая плотность $\text{PSD}\approx3.5\times10^{-4}$ при числе атомов в МОЛ порядка $N \approx 2\times 10^6$ и температуре $T\approx1\unit{мкК}$
\end{itemize}